\chapter{Análisis de la aplicación}
\label{cap:analisis_aplicacion} 

\section{Especificación de requisitos}
\label{sec:requisitos}
La especificación de requisitos es una etapa clave en el desarrollo de un proyecto, ya que en ella se definen las funcionalidades y condiciones que debe cumplir el sistema. Su objetivo principal
es establecer de manera clara y estructurada las expectativas del producto final.

En esta sección se identifican y detallan los requisitos, tomando en cuenta las necesidades y restricciones detalladas por los 
usuarios y demás partes interesadas. Para ello, es fundamental realizar un análisis que permita comprender qué características son esenciales 
para el correcto funcionamiento del software.

Cada requisito contará con un código único que permitirá clasificarlo y referenciarlo fácilmente a lo largo del documento. Este método de organización facilita su seguimiento y garantiza que 
todas las especificaciones sean consideradas durante el desarrollo del proyecto. 

\subsection{Requisitos funcionales}
\label{subsec:requisitos_funcionales}
Los requisitos funcionales son especificaciones detalladas que describen las funciones y comportamientos que un sistema debe cumplir para 
satisfacer las necesidades del usuario. Estos requisitos definen las acciones que el sistema debe ser capaz de realizar, incluyendo las interacciones
con el usuario y las respuestas a entradas específicas. Estos requisitos son ensenciales para guiar el desarrollo del sistema y asegurar que cumpla con las expectivas del usuario. \cite{pmo2017}

Vamos a clasificar estos requisitos en diferentes grupos según su funcionalidad. 

\begin{itemize}
    \item \textbf{Grupo 1:} Gestión de sesión. En este grupo se incluyen los requisitos relacionados con la autenticación y la gestión de la sesión del usuario.
    \item \textbf{Grupo 2:} Gestión del hogar. En este grupo se incluyen los requisitos relacionados con la creación y gestión de un hogar compartido.
    \item \textbf{Grupo 3:} Gestión de tareas. En este grupo se incluyen los requisitos relacionados con la creación y asignación de tareas en el hogar.
    \item \textbf{Grupo 4:} Gestión de gastos. En este grupo se incluyen los requisitos relacionados con la creación y gestión de gastos en el hogar.
    \item \textbf{Grupo 5:} Gestión de compra. En este grupo se incluyen los requisitos relacionados con la creación y gestión de listas de compra.
\end{itemize}

A continuación, se detallan los requisitos funcionales de cada grupo.

\subsubsection{Grupo 1: Gestión de sesión}
\label{subsubsec:grupo1}


\begin{table}[htb!]
	\centering 
	\begin{tabular}{|p{0.35\linewidth}|p{0.6\linewidth}|}
		\hline
		\rowcolor[HTML]{C0C0C0}  
        \textbf{RF\_01} & \textbf{Registro} \\
		\hline
		\textbf{Prioridad} & Alta \\
		\hline
		\textbf{Descripción} & El sistema debe permitir al usuario registrarse.  \\
		\hline
		\vspace{0.5mm}
		\textbf{Criterios de aceptación} & 
		\begin{minipage}[t]{0.9\linewidth}
			\begin{enumerate}
				\item La interfaz de registro debe contener los campos: nombre, apellidos, correo electrónico y contraseña.
				\item Verificar que los campos no estén vacíos.
				\item Verificar que la información introducida sea válida y que los campos cumplan con el formato requerido.
				\item Si los campos no son válidos o ya pertenecen a un usuario registrado, mostrar un mensaje de error.
				\item Tras el registro, se mostrará un mensaje indicando que debe verificar su correo electrónico y se redigirá a la pantalla de inicio de sesión.
			\end{enumerate}
			\vspace{2mm}
		\end{minipage} \\
		\hline
	\end{tabular}
	\caption{RF\_01}
\end{table}

\begin{table}[htb!]
    \begin{tabular}{|p{0.35\linewidth}|p{0.6\linewidth}|}
        \hline
        \rowcolor[HTML]{C0C0C0}  
        \textbf{RF\_02} & \textbf{Inicio de sesión} \\
        \hline
        \textbf{Prioridad} & Alta \\
        \hline
        \textbf{Descripción} & El sistema debe permitir al usuario iniciar sesión.  \\
        \hline
        \vspace{0.5mm}
        \textbf{Criterios de aceptación} & 
        \begin{minipage}[t]{0.9\linewidth}
            \begin{enumerate}
                \item La interfaz de inicio de sesión debe contener los campos: correo electrónico y contraseña.
                \item Verificar que los campos no estén vacíos.
                \item Verificar que la información introducida sea válida y que los campos cumplan con el formato requerido.
                \item Si los campos no son válidos o no pertenecen a un usuario registrado, mostrar un mensaje de error.
                \item Si el usuario no ha verificado su correo electrónico, mostrar un mensaje indicando que debe verificarlo.
                \item Tras el inicio de sesión, se redirigirá a la pantalla principal de la aplicación.
            \end{enumerate}
            \vspace{2mm}
        \end{minipage} \\
        \hline
    \end{tabular}
    \caption{RF\_02}

\end{table}


\begin{table}[htb!]
    \begin{tabular}{|p{0.35\linewidth}|p{0.6\linewidth}|}
        \hline
        \rowcolor[HTML]{C0C0C0}  
        \textbf{RF\_03} & \textbf{Inicio de sesión con Google} \\
        \hline
        \textbf{Prioridad} & Baja \\
        \hline
        \textbf{Descripción} & El sistema debe permitir al usuario iniciar sesión con Google  \\
        \hline
        \vspace{0.5mm}
        \textbf{Criterios de aceptación} & 
        \begin{minipage}[t]{0.9\linewidth}
            \begin{enumerate}
                \item La interfaz de inicio de sesión debe contener un botón para iniciar sesión con Google.
                \item Al pulsar el botón, se abrirá una ventana emergente para iniciar sesión con la cuenta de Google.
                \item Si el usuario ya ha iniciado sesión con Google, se redirigirá a la pantalla principal de la aplicación.
                \item Si el usuario no ha iniciado sesión con Google, se registrará automáticamente y se redirigirá a la pantalla principal de la aplicación.
            \end{enumerate}
            \vspace{2mm}
        \end{minipage} \\
        \hline
    \end{tabular}
    \caption{RF\_03}
\end{table}

\begin{table} [htb!]
    \begin{tabular}{|p{0.35\linewidth}|p{0.6\linewidth}|}
        \hline
        \rowcolor[HTML]{C0C0C0}  
        \textbf{RF\_04} & \textbf{Cerrar Sesión} \\
        \hline
        \textbf{Prioridad} & Alta \\
        \hline
        \textbf{Descripción} & El sistema debe permitir al usuario cerrar sesión  \\
        \hline
        \vspace{0.5mm}
        \textbf{Criterios de aceptación} & 
        \begin{minipage}[t]{0.9\linewidth}
            \begin{enumerate}
                \item La interfaz de la aplicación debe contener un botón para cerrar sesión.
                \item Al pulsar el botón, se mostrará un mensaje de confirmación.
                \item Si el usuario confirma, se cerrará la sesión y se redirigirá a la pantalla de inicio de sesión.
            \end{enumerate}
            \vspace{2mm}
        \end{minipage} \\
        \hline
    \end{tabular}
    \caption{RF\_04}

\end{table}








